\section{Conclusion}
\label{sec:conclusion}

We studied a simple but consequential question: whether death is a uniquely disqualifying topic for LLM writing. Using 12 human grief passages, 24 matched frontier-model generations, blinded LLM judges, a pretrained detector, and lightweight style analysis, we find a small human-model authenticity gap that does not grow for bereavement. The main interaction term is near zero, while detector separability remains substantial.

The key takeaway is straightforward. In this automated pilot, lack of firsthand experience did not make LLM death-writing categorically unlike human grief writing. Generated texts remained mechanically distinguishable, but that distinguishability did not translate into strong authenticity penalties from judges.

The next steps are clear. Future work should replace LLM judges with blinded human raters, expand to a larger and more naturalistic bereavement corpus, compare process-authenticity with text-authenticity directly, and use better-calibrated rubrics that avoid ceiling effects. Those extensions would sharpen the boundary between what readers perceive, what detectors measure, and what philosophers mean when they call a piece of writing authentic.
